\section*{Zakończenie}\section*{Zakończenie}
\addcontentsline{toc}{section}{Zakończenie}
\label{sec:zakonczenie}
% dotad sprawdzono

Niniejsza praca magisterska była poświęcona wielokryterialnej analizie i ocenie przydatności współczesnych technologii automatycznego rozpoznawania mowy (\textit{Speech-to-Text}, STT) w kontekście wsadowego przetwarzania nagrań audio dla języka polskiego. Realizacja celu głównego wymagała połączenia studiów literaturowych z obszaru cyfrowego przetwarzania sygnałów i głębokiego uczenia maszynowego, prac programistycznych nad autorskim stanowiskiem badawczym oraz przeprowadzenia cyklu pomiarów empirycznych na zróżnicowanym materiale audio.

\subsection*{Podsumowanie pracy}

W ramach przeprowadzonego wywodu naukowego i inżynieryjnego zrealizowano wszystkie założone cele szczegółowe:
\begin{enumerate}
    \item \textbf{Syntetyczna podstawa teoretyczna:} Przedstawiono ewolucję paradygmatów ASR -- od klasycznych Ukrytych Modeli Markowa (HMM) i hybryd HMM-DNN, przez funkcję celu CTC, po zintegrowane architektury sekwencja-do-sekwencji (\textit{Transformer}, \textit{Conformer}) oraz uczenie samo-nadzorowane (\textit{wav2vec 2.0}) i słabo nadzorowane na wielką skalę (\textit{OpenAI Whisper}).
    \item \textbf{Przegląd rozwiązań rynkowych:} Przeprowadzono szczegółową klasyfikację i analizę porównawczą komercyjnych usług chmurowych (Google Cloud STT v2, Microsoft Azure Speech Service, Amazon Transcribe, OpenAI API), otwartych modeli głębokich (Whisper, WhisperX) oraz natywnych interfejsów przeglądarkowych (Web Speech API).
    \item \textbf{Autorskie stanowisko badawcze:} Zaprojektowano i zaimplementowano aplikację webową w architekturze kliencko-serwerowej (FastAPI + React 19 + TypeScript + Vite), zoptymalizowaną pod kątem wsadowego przetwarzania nagrań audio, asynchronicznej orkiestracji wywołań interfejsów API, trwałej archiwizacji wyników w plikach JSON oraz wizualizacji alokacji błędów (diff-HTML).
    \item \textbf{Ewaluacja empiryczna i weryfikacja hipotezy:} Przeprowadzono badanie na zbiorze obejmującym około 300 transkrypcji zrealizowanych przez grupę 10 probantów. Sformułowaną w pracy hipotezę badawczą uznano za \textbf{częściowo potwierdzoną}:
    \begin{itemize}
        \item W wymiarze dokładności transkrypcji (WER/CER) hipoteza potwierdziła się dla zoptymalizowanego wariantu \textbf{WhisperX}, który osiągnął najniższy średni wskaźnik $\text{WER} \approx 11,4\%$ oraz $\text{CER} \approx 4,6\%$. Wskazano jednak, że natywny model \textit{Whisper Offline} bez wstępnej detekcji VAD cierpiał na zjawisko halucynacji i pętlenia na ciszy, ustępując wskaźnikami niektórym usługom chmurowym.
        \item W wymiarze wydajności czasowej (RTF) hipoteza potwierdziła się w pełni -- komercyjne interfejsy API (\textit{OpenAI GPT-4o Transcribe} z $\text{RTF} = 0,128$ oraz \textit{Google STT} z $\text{RTF} = 0,211$) wykazały wielokrotnie wyższą wydajność czasową niż lokalne modele z rodziny Whisper ($\text{RTF} \approx 1,62$).
    \end{itemize}
\end{enumerate}

\subsection*{Najważniejsze osiągnięcia}

Do kluczowych osiągnięć naukowych i practicalnych niniejszej pracy należą:
\begin{itemize}
    \item \textbf{Stworzenie powtarzalnego środowiska ewaluacyjnego:} Opracowanie bezstanowej, asynchronicznej aplikacji badawczej z odizolowanymi adapterami STT, pozwalającej na obiektywne porównywanie silników ASR bez ryzyka blokowania interfejsu czy utraty wyników cząstkowych.
    \item \textbf{Wykazanie kluczowej roli modułów VAD i forced alignment:} Dowiedzenie na materiale empirycznym, że wstępna segmentacja mowy oraz phoneme-level forced alignment (zaimplementowane w WhisperX) są warunkiem koniecznym do wyeliminowania halucynacji oraz zapętlania tekstu w modelach autoregresywnych.
    \item \textbf{Identyfikacja ograniczeń funkcji ITN w usługach chmurowych:} Wykazanie, że nieujednolicony format zapisu danych liczbowych i godzin (zjawisko \textit{Inverse Text Normalization}) w usługach Azure i Amazon generuje sztuczne skoki wskaźnika WER, co uzasadnia konieczność stosowania warstw post-normalizacji tekstu w aplikacjach produkcyjnych.
    \item \textbf{Formułowanie rekomendacji architektonicznych:} Opracowanie wytycznych wdrożeniowych dzielących zastosowania na systemy interaktywne o niskiej latencji (rekomendacja: GPT-4o / Google STT v2) oraz aplikacje wsadowe o najwyższej dokładności i wymogach prywatności (rekomendacja: WhisperX / Self-Hosted GPU).
\end{itemize}

\subsection*{Kierunki dalszych badań}

Obszar automatycznego rozpoznawania mowy rozwija się niezwykle dynamicznie. Wśród najbardziej obiecujących kierunków dalszych prac naukowych i programistycznych wyróżnić można:
\begin{enumerate}
    \item \textbf{Integracja z wielomodalnymi modelami językowymi (Speech LLMs):} Przejście od klasycznego kaskadowego przetwarzania ASR$\to$NLU do bezpośredniego dekodowania tokenów audio przez modele wielomodalne (\textit{Native Speech-to-Speech}), co pozwala na zachowanie informacji o intonacji, emocjach oraz tożsamości mówcy.
    \item \textbf{Inference po stronie klienta (Edge ASR / WebAssembly):} Kompilacja lekkich modeli neuronowych do formatów ONNX Web / WASM, co umożliwi wykonywanie pełnej transkrypcji bezpośrednio w przeglądarce użytkownika bez konieczności angażowania serwerowych zasobów GPU.
    \item \textbf{Semantyczne metryki oceny jakości:} Zastąpienie lub uzupełnienie tradycyjnych wskaźników Levenshteina (WER/CER) metrykami opartymi na zanurzeniach pojęciowych (np. \textit{BERTScore} lub ocena semantyczna z użyciem LLM), mierzącymi stopień zachowania rzeczywistego sensu wypowiedzi.
    \item \textbf{Hybrydowy routing zapytań STT:} Zaimplementowanie w warstwie backendu brokerów zapytań, które na podstawie wstępnej analizy czasu trwania, poziomu szumu oraz wymogów prywatności dynamicznie kierują strumień audio do najbardziej optymalnego silnika ASR.
\end{enumerate}

\medskip

Reasumując, przeprowadzone badania empiryczne oraz zaimplementowane środowisko ewaluacyjne dowodzą, że technologia automatycznego rozpoznawania mowy w języku polskim osiągnęła poziom dojrzałości umożliwiający szerokie zastosowanie produkcyjne. Wybór optymalnego rozwiązania nie jest jednak uniwersalny i wymaga precyzyjnego balansu pomiędzy dokładnością transkrypcji, wydajnością czasową a architekturą systemu. Opracowane w niniejszej pracy wnioski oraz wytyczne stanowią praktyczny przewodnik dla inżynierów i architektów projektujących nowoczesne systemy wsadowego przetwarzania nagrań audio.